\section{审判制度与法官职业道德}
\begin{enumerate}
  \item 法检平行，要求一样。
  \item 审判制度。
        \begin{enumerate}
          \item 审判原则。司法公正、审判独立、不告不理、直接言词、及时审判。
          \item 审判机关。巡回法院 = 最高法本部。军事法院 3 级。专门法院级别中级。\zd{设立法院依据。宪法、全人常}。
          \item 审判组织。独任庭、合议庭、人民陪审员、审委会（司法解释）、专业委员会（指导案例）。
          \item \kd{法官任免}。检察也同。除了院长可以不过法考。
                \begin{enumerate}
                  \item 至少本科。\zd{本硕博，543}。
                  \item \zd{案底绝对干净}。
                  \item 限制性条件。
                        \begin{enumerate}
                          \item 兼职。\zd{学术性兼职 + 单位同意}。
                          \item 离职 2 年以内、律师身份、诉讼业务。永远不要回来，近亲辩护是例外。
                          \item 任职回避。同院不同庭。\yi（副）庭长 + 审判员。\er 都是审判员。不同院\zd{相邻两级正副职领导}。
                        \end{enumerate}
                  \item \kd{遴选}。最高和高院遴选委员会 1/3 自己人。惩戒委员会 1/2。
                        \begin{enumerate}
                          \item 初任法官，下方基层。一般逐级调。
                          \item 院长、副院长要求法律经验。
                          \item 扩大范围。律师、科研人员 + 5 年执业/科研经历。
                          \item 高院全人大，其他同级人常。
                        \end{enumerate}
                  \item 处分、考评等同公务员。
                  \item 退休保障。
                  \item 权利义务。\zd{普法义务}。
                \end{enumerate}
        \end{enumerate}
  \item 法官职业道德。做题\zd{合理怀疑}。\zd{法官（及家属）不是人，不下班，不退休}。
        \begin{enumerate}
          \item 对象。扩大到人民陪审员、退休法官。
          \item 内容。忠诚、\zd{为民、公正、廉洁、法官讲形象，检察官有担当}。
          \item \zd{不过问、不干预、不评论}。题干、选项涉及在审案件基本都是错的。
          \item 职业责任。
          \item \zd{不予处分。没问题}。对法律、证据认识偏差。
          \item \zd{免予处分。有问题但很小。情节轻，无后果，态度好，能改正}。
          \item 解除。
          \item 变更 \hz 撤销。处分不当。
        \end{enumerate}
  \item 人民陪审员。参照刑诉。
        \begin{enumerate}
          \item 排除公职、专业限制。要的是常识。
          \item 5 倍海选，3 倍正式，人常任命，推荐最多 1/5。
          \item 5 年，不得连任。
          \item 案情复杂，影响大。
        \end{enumerate}
\end{enumerate}
